\documentclass[11pt]{article}

\usepackage[margin=1in]{geometry}
\usepackage{amsmath,amssymb,amsthm,mathtools}
\usepackage{booktabs}
\usepackage{enumitem}
\usepackage{tabularx}
\usepackage{hyperref}
\usepackage{microtype}
\usepackage[nameinlink,noabbrev]{cleveref}

\newtheorem{theorem}{Theorem}
\theoremstyle{definition}
\newtheorem{definition}{Definition}
\theoremstyle{remark}
\newtheorem{remark}{Remark}

\newcommand{\E}{\mathbb{E}}

\title{State-Dependent Anonymity (Addendum):\\Calibration Vignettes and Time-to-Identification Conversions}
\author{}
\date{March 19, 2026 (archive v393)}

\begin{document}
\maketitle

\begin{abstract}
This companion note collects only the small calibration-support packet for the State-Dependent Anonymity model (State~1).
It is intentionally non-redundant: we do not restate the full game definitions, amplification bounds, or public release machinery.
Instead we record two reusable translations:
(i) converting a per-epoch witness gap into time-to-identification, and
(ii) converting a per-lookup MaxL bound into a Boss-Fight horizon budget.
The note is intentionally narrow: it provides referee-quotable conversions and vignette tables, not a second theorem root, evaluator notarization surface, or release-packaging layer.
\end{abstract}

\paragraph{Scope.}
Calibration addendum for State~1.\cite{series:state1}

\paragraph{Units.} Budgets are published in bits (use $\log_2$ and $2^{\epsilon}$); results stated in nats can be converted by dividing by $\ln 2$ (see Synthesis~8).\cite{series:notation}

\paragraph{Imports.}
We import the state-dependent anonymity viewpoint and witness taxonomy from State~1, along with the endpoint sizing vocabulary and the repeated-use horizon accountant from the endpoint bridge and Boss Fight~A.\cite{series:state1,series:endpointbridge,series:bossfightA}
The surrounding evaluation / publication stack is imported rather than re-owned here: evaluator notarization belongs to Eval~1, and public claim packaging / release lineage belong to ABOM/OINL plus Release~A.\cite{series:evalnotary,series:abomoinl,series:releaseproperty}
Artifact availability follows the artifact policy.\cite{series:artifactpolicy}

\paragraph{Exports.}
A compact calibration-support packet: two reusable conversions, one small epoch-to-time table, and one minimal declaration pattern that help a referee interpret ``small'' per-epoch gaps without reopening the full theorem or release stack.

\paragraph{Boundary.}
This addendum owns only those supplementary conversions and vignette tables.
It does \emph{not} own the state-anonymity games / accountant root (State~1), the repeated-use horizon accountant (Boss Fight~A), evaluator notarization / anti-grinding evidence release (Eval~1), or public claim packaging / release lineage (ABOM/OINL plus Release~A).\cite{series:state1,series:bossfightA,series:evalnotary,series:abomoinl,series:releaseproperty}

\paragraph{Exported object.}
The exported object is a small calibration packet: one gap-to-time rule of thumb, one MaxL-to-horizon conversion rule, one vignette table translating repeated public witnesses into time scales, and one minimal declaration pattern that tells a reader which cadence / horizon assumptions must accompany the conversion.

\paragraph{Earliest sufficient stop.}
This note is complete once a reader can translate a declared State~1 witness gap or per-lookup MaxL budget into a repeated-use time or query horizon and can cite the corresponding vignette table.
If a deployment needs the underlying theorem packet, a replayable evaluator log, or a public release-lineage bundle, the reader should leave this addendum and use State~1, Eval~1, and ABOM/OINL + Release~A directly.\cite{series:state1,series:evalnotary,series:abomoinl,series:releaseproperty}

\paragraph{Later terse-paper citation posture.}
For later terse papers under space pressure, default to State~1 for the state-anonymity theorem root and to Boss Fight~A for repeated-use horizon claims; cite this addendum only when the exact calibration-support packet exported here --- the gap-to-time rule, the MaxL-to-horizon conversion, the vignette table, or the minimal declaration pattern --- is itself live.
If the live question is only the public theorem/accountant claim, stop at those main notes.
If the live question is instead ``how was the quoted time-scale or horizon translation obtained?'', widen to this addendum.
If the question becomes evaluation notarization or public release-lineage comparison, leave this addendum and reopen the neighboring owners.
The maintained worked example and paired cutover rule are the general model for that stop / widen / reopen discipline.\cite{series:workedexample,series:pairedterminalcutoverrules}

\paragraph{Stable handoff to neighboring layers.}
This addendum is easiest to audit when the stopping rule is stated directly rather than inferred from surrounding series context.
The handoff is therefore deliberately short and explicit.

\begin{center}
\begin{tabularx}{0.97\linewidth}{@{}>{\raggedright\arraybackslash}p{0.31\linewidth}>{\raggedright\arraybackslash}p{0.34\linewidth}>{\raggedright\arraybackslash}X@{}}
\toprule
\textbf{Referee question} & \textbf{Earliest sufficient stop} & \textbf{Owner} \\
\midrule
How do I convert a declared state witness gap into a time scale I can quote? & This addendum's calibration-support packet: gap-to-time rule, vignette table, and declared witness cadence. & State~1A \\
What theorem/accountant justifies the state-dependent witness model in the first place? & State-anonymity game, witness taxonomy, and theorem/accountant root. & State~1 \\
How do repeated lookups accumulate into a horizon budget? & Repeated-use accountant and horizon-spend discipline. & Boss Fight~A \\
Was the claimed deployment actually evaluated and logged under a notarized discipline? & Evaluator manifest, attempt log, and anti-grinding receipt. & Eval~1 \\
How should the final public claim be packaged and compared to earlier releases? & Lineaged release receipt / release-ledger entry. & ABOM/OINL plus Release~A \\
\bottomrule
\end{tabularx}
\end{center}


That table is intentionally narrow: the addendum exports calibration support for State~1 claims, not a competing theorem, repeated-use accountant, evaluator receipt, or publication object.

\paragraph{A minimal public calibration-support card.}
When a later short paper only needs to answer \emph{which calibration packet turned this declared state witness or per-lookup MaxL bound into a quotable time / horizon statement?}, it should stop at the small card below plus the tiny drill that follows.

\begin{table}[t]
\centering
\begin{tabularx}{0.97\linewidth}{@{}>{\raggedright\arraybackslash}p{0.31\linewidth}>{\raggedright\arraybackslash}X@{}}
\toprule
\textbf{Field} & \textbf{Declared calibration support object} \\
\midrule
State theorem anchor & \texttt{state\_claim\_id = state1:witness-gap-v1}; ties the calibration packet back to the State~1 witness model rather than creating a second theorem root. \\
Witness cadence & \texttt{cadence = 1 witness / hour}; the public cadence assumption that turns epoch counts into wall-clock time. \\
Gap-to-time rule & \texttt{T \textasciitilde 2 log2(1/eta) / Delta\^{}2}; the reusable rule-of-thumb from \Cref{thm:repeat-witness}. \\
Target error & \texttt{eta = 0.05}; fixes the horizon row a referee is allowed to quote from the vignette table. \\
Vignette row & \texttt{Delta = 0.02 -> T \textasciitilde 1.5e4 epochs -> 1.7 years @ 1 hour}; one directly quotable conversion row. \\
Per-lookup MaxL line & \texttt{epsilon = 2e-4 bits / lookup}; the declared per-lookup MaxL budget for the same projection. \\
Repeated-use horizon line & \texttt{B = 0.2 bits -> Q <= 1000 lookups}; the conservative Boss-Fight carry-forward derived from the declared per-lookup bound. \\
\bottomrule
\end{tabularx}
\caption{A minimal public calibration-support card. This is the smallest public answer later terse papers can quote directly when the live issue is how a declared State~1 witness gap or MaxL lookup cap was converted into a time or horizon statement.}
\label{tab:state1a-card}
\end{table}

\paragraph{Tiny calibration drill.}
Suppose a deployment declares a witness gap $\Delta=0.02$, target total error $\eta=0.05$, and one public witness per hour.
\Cref{thm:repeat-witness} then gives $T \approx 1.5\times 10^{4}$ epochs, which is about $1.7$ years at that cadence; the same packet can declare a per-lookup MaxL cap $\varepsilon=2\times 10^{-4}$ bits, so a conservative repeated-use budget $B=0.2$ bits permits at most
\[
Q \le \frac{B}{\varepsilon} = \frac{0.2}{2\times 10^{-4}} = 1000
\]
lookups before the note must widen to Boss Fight~A for a fresh horizon statement.
That is enough for a later terse paper to quote one calibration packet without reopening the theorem root, evaluator receipts, or release lineage.

\section{Gap-to-time conversion for repeated public witnesses}

\begin{definition}[Per-epoch witness gap]
Let $W_t\in\{0,1\}$ be a public witness bit observed once per epoch (reconfiguration, consensus, or lookup),
with $\Pr[W_t=1\mid C=1]=p_1$ and $\Pr[W_t=1\mid C=0]=p_0$.
Define the per-epoch gap $\Delta := |p_1-p_0|$.
\end{definition}

\begin{theorem}[Rule-of-thumb identification horizon]\label{thm:repeat-witness}
Fix target total error $\eta\in(0,1)$.
For i.i.d.\ Bernoulli witnesses with gap $\Delta$, a simple likelihood-ratio test achieves error $\le \eta$ once
\[
T \;\gtrsim\; \frac{2\log_2(1/\eta)}{\Delta^2}.
\]
Conversely, any test requires $T=\Omega(\log_2(1/\eta)/\Delta^2)$ epochs when $\Delta$ is small.
\end{theorem}

\paragraph{Epoch-to-time.}
To translate epochs into wall-clock time, pick a witness cadence (per consensus period, per reconfiguration, or per lookup).
Table~\ref{tab:vignette-times} instantiates \Cref{thm:repeat-witness} at $\eta=0.05$.

\begin{table}[t]
\centering
\begin{tabular}{@{}r r r r@{}}
\toprule
Gap $\Delta$ & Required epochs $T$ & Time @ 5 min & Time @ 1 hour \\
\midrule
0.01 & $6.0\times 10^{4}$ & $\approx 208$ days & $\approx 6.8$ years \\
0.02 & $1.5\times 10^{4}$ & $\approx 52$ days & $\approx 1.7$ years \\
0.05 & $2.4\times 10^{3}$ & $\approx 8.3$ days & $\approx 100$ days \\
\bottomrule
\end{tabular}
\caption{Epoch-to-time conversion for \Cref{thm:repeat-witness} at $\eta=0.05$.}
\label{tab:vignette-times}
\end{table}

\paragraph{When epochs are ``per lookup.''}
If the witness occurs once per lookup and a client issues $r$ lookups/day, then $T$ epochs corresponds to $T/r$ days.
This is one reason receipts often publish a horizon budget (queries per window) rather than a clock-time guarantee.

\section{Per-lookup MaxL to Boss-Fight horizon}

If a deployment declares a per-lookup MaxL bound of the form
\[
\mathcal{L}_{\max}(C\!\to\! Y)\le \varepsilon
\]
for the declared projection $Y$, then under the conservative repeated-use accountant
\[
\mathcal{L}_{\max}(C\!\to\! Y_{1:Q}) \le Q\,\varepsilon.
\]
This bound is tight in adversarial stop/abort settings unless additional independence or reset structure is declared.\cite{series:bossfightA}

\paragraph{Minimal declaration line item.}
Declare $(Q,\varepsilon)$ (or a tiered list $(Q_i,\varepsilon_i)$) together with the witness cadence used to interpret $Q$.
If that declaration is later exported publicly, its packaging belongs to the receipt schema / ABOM-OINL / Release~A stack rather than to this calibration note.\cite{series:receiptitems,series:abomoinl,series:releaseproperty}
That declaration line should therefore be read as a small calibration companion to the State~1 theorem root and Boss Fight~A accountant rather than as a shadow receipt or release template.

\begin{thebibliography}{99}

\bibitem{series:state1}
Anonymous.
\newblock State-Dependent Anonymity: Selection and State Evolution as an Attack Surface.
\newblock AnonDHT State series Paper~1, 2026. (This archive.)

\bibitem{series:bossfightA}
Anonymous.
\newblock Boss Fight Budgets: Repeated-Use Identification Accounting for Anonymous DHT Lookups.
\newblock Boss Fight series Paper~A, 2026. (This archive.)

\bibitem{series:evalnotary}
Anonymous.
\newblock Notarized Anonymity Certificates: Evaluator Receipts, Spending Discipline, and Anti-Grinding Evidence.
\newblock Evaluation series Paper~1, 2026. (This archive.)

\bibitem{series:abomoinl}
Anonymous.
\newblock ABOM/OINL: Audit BOM and Outcome Identification Noise Label for Anonymous DHT Deployments.
\newblock Anonymity series Paper~C, 2026. (This archive.)

\bibitem{series:releaseproperty}
Anonymous.
\newblock Release Property Ledgers and Receipts for Anonymous DHT Privacy Claims.
\newblock Release and Destination series Paper~A, 2026. (This archive.)

\bibitem{series:endpointbridge}
Anonymous.
\newblock Endpoint Metrics Bridge: Odds Inflation, PML/MaxL, and Identification Sizing.
\newblock Series synthesis note (Synthesis~5), 2026. (This archive.)

\bibitem{series:receiptitems}
Anonymous.
\newblock Receipt Line-Item Schema for Anonymous-DHT Privacy Claims.
\newblock Series synthesis note (Synthesis~12), 2026. (This archive.)

\bibitem{series:notation}
Anonymous.
\newblock Notation and Minimal Model for the One-True-Archive.
\newblock Synthesis~8, The-One-True-Archive (2026).

\bibitem{series:workedexample}
Anonymous.
\newblock Worked Example: Instantiating a Tiered Reader-Privacy Receipt.
\newblock Synthesis~17, The-One-True-Archive (2026).

\bibitem{series:pairedterminalcutoverrules}
Anonymous.
\newblock Paired Terminal Cutover Rules for Stable Review Spines: When Later Terse Papers Must Stop, Widen, or Reopen.
\newblock Series synthesis note (Synthesis~75), 2026. (This archive.)

\bibitem{series:artifactpolicy}
Anonymous.
\newblock Artifact policy and supporting materials for the anonymity/AnonDHT series.
\newblock 2026.
\end{thebibliography}

\end{document}
